\chapter{Теория}

\section{Состав и частичная расшифровка MAC-адреса}
MAC-адрес состоит их 6 октетов, записанных в виде 00:0c:29:1c:d1:fa. Такая запись используется в ОС Ubuntu, но другие ОС могут записывать MAC-адрес в другой форме (00-0C- 29-AD-CE-F2; 000c.29ad.cef2).
Разберем на примере состав MAC-адреса:

\begin{figure}[H]
    \centering
    \includegraphics[width=0.75\textwidth]{mac.png}
    \caption{Расшифровка MAC адреса}
\end{figure}

Первые три байта – OUI – уникальный номер, который указывает на организацию, которая изготовила данную сетевую карту. \\

Вторые три байта – NIC Specific – это оригинальный идентификатор конкретного устройства.
Эта часть задается производителем уникально для каждого устройства. \\

После определения MAC-адреса возможно узнать организацию, которая изготовила вашу сетевую карту. Для этого можно воспользоваться
определенными сайтами (их много, они легко ищутся в поисковике).


\endinput